\documentclass[a4paper,11pt]{article}
\usepackage[utf8]{inputenc}
\usepackage[norsk]{babel}
\usepackage{amsmath}
\usepackage{graphicx}
\usepackage{geometry}
\usepackage{fancyhdr}
\usepackage{booktabs}
\usepackage{array}
\geometry{margin=2.5cm}
\pagestyle{fancy}
\fancyhf{}
\lhead{Laboratorieoppgave: Reguleringsteknikk}
\rhead{Prosjekt: Den Digitale Tvillingen}
\cfoot{\thepage}

\begin{document}

\section*{Lab-oppgave: Systemidentifikasjon og SIMC-tuning}

\textbf{Navn:} \underline{\hspace{5cm}} \textbf{Dato:} \underline{\hspace{3cm}} \textbf{Anlegg nr:} \underline{\hspace{1cm}}

\subsection*{1. Identifikasjon av prosessen blå (dataene dine) og hvit (matematisk modell) kurve.}
\textit{Bruk piltastene til å tilpasse den hvite modellen til den blå målingen din.}
\begin{itemize}
 \item Dødtid ($L$): \underline{\hspace{2cm}} sekunder
 \item Tidskonstant ($T$): \underline{\hspace{2cm}} sekunder
    \item Endring i nivå ($dY$): \underline{\hspace{2cm}}\% 
    \item Startverdei nivå ($Y_0$): \underline{\hspace{2cm}}\%
\end{itemize}
Så kan du lagre et bilde av grafen (trykk på lagre bilde, og bruk navnet ditt som filnavn (f.eks. marcus\_1).  Gi beskjed til lærer så overfører vi fila til minnepinne.
\subsection*{2. Velg $\lambda$ verdi: }
\textit{Når vi har en regulering så ønsker vi ofte å gjøre den så rask som mulig, samtidig som vi ønsker at systemer er stabilt. Dette justerer vi med å velge $\lambda$ verdi.  I programet kan vi velge mellom følgende $\lambda$ verdier:  }\\
\begin{table}[htbp]
\centering
\renewcommand{\arraystretch}{1.3} 
\begin{tabular}{|l|c|c|>{\centering\arraybackslash}p{2.5cm}|c|c|} % K_P er nå 2.5cm bred
\hline
Valgt $\lambda$ & Fartsfaktor ($T/\lambda$) & Integraltid ($T_i$) & $K_P$ & Oversving (\%) & Stabil? (J/N) \\ \hline
$\lambda = T$        & 1                         &                     &       &                &               \\ \hline
$\lambda = T/2$      & 2                         &                     &       &                &               \\ \hline
$\lambda = T/4$      & 4                         &                     &       &                &               \\ \hline
$\lambda = T/6$      & 6                         &                     &       &                &               \\ \hline
$\lambda = T/8$      & 8                         &                     &       &                &               \\ \hline
Egendefinert: \dots  &                           &                     &       &                &               \\ \hline
\end{tabular}
\end{table} \\
a) Hvordan vil vil valget av $\lambda$ påvirke reguleringen med tanke på hastighet og stabilitet? Hint: uten regulering er $\lambda=T$\\\\
b) Velg $\lambda$=T/4 og $\lambda=T/8$ og skriv inn dataene i tabellen over for den $\lambda$ verdien du valgte.\\\\
c) Fullfør en sprangrespons i set-punkt (la tanken reguleres tilbake til utgangspunktet først).  Så kan du lagre et bilde av grafen (trykk på lagre bilde, og bruk navnet ditt som filnavn (f.eks. marcus\_2).  Gi beskjed til lærer så overfører vi fila til minnepinne.
\subsection*{3. Analyse og diskusjon}
\begin{enumerate}
    \item Hva skjer med nivåprosessen når $\lambda$ blir liten?  Hva skjer med $K_P$ verdien da? 
    \item Studer kurven du har fått og vurder resultatet. Stikkord: stabilitet, hurtighet og oversving.  
    \item Vurder om prosessen din har fått en høyere eller lavere  $\lambda$ verdi enn det du har valgt.
\end{enumerate}
\end{document}
